\chapter{E2: Multi-head self-attention và tokenization}
\section{Mục tiêu và biểu diễn token}
E2 hiện thực MSA bằng Linear, reshape, matmul, softmax và concat; bản đối
chiếu dùng \texttt{nn.MultiheadAttention}. Encoder và classifier giống nhau
giữa hai bản. Ba tokenizer cùng dùng split/normalization Fashion-MNIST ở
Chương 1; không sử dụng checkpoint E1.

\begin{figure}[H]\centering
\resizebox{.97\linewidth}{!}{\ArchTokens}
\caption{Ba tokenizer E2; CNN-stem có hai Conv kernel 3, stride 1, padding 1, bias, ReLU và MaxPool 2 sau mỗi Conv. Rows không có Conv; patch lấy trực tiếp pixel.}
\end{figure}
Rows có $ N=28,F=28$; patch/CNN-stem có $ N=49,F=16$.
CNN-stem thêm 2.480 tham số học so với patch cùng chiều feature, nên so sánh
tokenizer thay đổi cả inductive bias và capacity.

\section{Kiến trúc encoder và MSA}
\begin{figure}[H]\centering
\resizebox{.97\linewidth}{!}{\ArchEncoder}
\caption{Classifier E2: một FC embedding, hai encoder block và một FC head; không CLS token.}
\end{figure}
\begin{figure}[H]\centering
\resizebox{.97\linewidth}{!}{\ArchMSA}
\caption{Một pre-norm encoder block: hai projection MSA và hai FC trong FFN, tổng bốn ánh xạ Linear/block. Đường ngoài là residual.}
\end{figure}
Với $ D=64,H=4,d_h=16$, QKV reshape thành $ B\times H\times N\times d_h $:
\begin{equation}
A_h=\operatorname{softmax}_{\mathrm{key}}\left(Q_hK_h^\top/\sqrt{d_h}\right),
\qquad Y=\operatorname{Concat}_{h=1}^{H}(A_hV_h)W_O+b_O.
\end{equation}
Attention dropout 0,1 áp dụng lên $ A_h $ trước khi nhân $ V_h $; output dropout
nằm trên nhánh residual. Mỗi block có FFN $64\to 128\to 64$, GELU,
dropout 0,1. Toàn classifier có 10 ánh xạ Linear theo cách đếm QKV gộp là
một Linear: 1 embedding + $2\times4$ trong encoder + 1 head. CNN-stem bổ sung 2 Conv.

Cặp manual/PyTorch copy toàn bộ weights khởi tạo, gồm tokenizer, position,
FFN và head. Kiểm tra CUDA float64 với dropout tắt đạt sai số tối đa
$1,665\times10^{-16}$ ở forward và input/weight/bias gradients.
Parity xác nhận phép toán; không đảm bảo dropout/RNG và đường tối ưu của
hai lần training giống nhau.

\section{Thiết lập huấn luyện và tám cấu hình}
AdamW LR $10^{-3}$, WD $10^{-4}$, batch 128, tối đa 100 epoch; thí nghiệm dùng
patience 10, min-delta $10^{-4}$ và checkpoint min validation loss.
Sáu run cơ sở là manual/PyTorch $\times $ rows/patch/stem, H=4.
Tokenizer manual tốt nhất theo validation là CNN-stem, được dùng thêm H=2 và H=8;
H=4 tái dùng run đã train. D giữ 64 nên số tham số không thay đổi theo H.

Ký hiệu bảng: M=manual, T=PyTorch; R=rows, P=patch 4, C=CNN-stem;
chữ số cuối là H. Ví dụ M-C4 là manual CNN-stem, 4 head.

\section{Kết quả và so sánh theo đề}
\DataTable{assets/e2/results.tex}{Metric E2; mọi model đánh giá đủ 10.000 ảnh test, F1 là macro-F1.}
\DataTable{assets/e2/cost.tex}{Capacity và chi phí của tám cấu hình E2.}
\ReportFigure{assets/e2/comparison.png}{Tổng quan accuracy validation, capacity và chi phí E2.}
\ReportFigure{assets/e2/implementation_tokenizer_comparison.png}{So sánh manual/PyTorch cùng tokenizer: accuracy, params và giây/epoch.}
\begin{table}[H]\centering\caption{So sánh manual/PyTorch ở H=4, cùng weights khởi tạo.}
\small\begin{tabular}{lrrrr}\toprule Tokenizer & M val & T val & M--T (điểm\%) & M/T s/epoch\\\midrule
Rows &89,97&89,90&+0,07&6,217 / 6,142\\
Patch &89,17&88,47&+0,70&7,347 / 6,902\\
CNN-stem &90,65&90,68&$-0,03$&6,487 / 6,824\\\bottomrule\end{tabular}\end{table}
Sai số đối chiếu phép toán rất nhỏ không có nghĩa metric training phải bằng nhau:
dropout, thứ tự sử dụng số ngẫu nhiên và AMP làm đường tối ưu khác.
Chênh lệch 0,07 hoặc 0,03 điểm ở rows/stem nhỏ; chênh lệch 0,70 ở patch cũng
chỉ quan sát ở một seed, chưa kết luận implementation nào tốt hơn.

CNN-stem đạt validation cao nhất trong cả hai implementation
(90,65\% manual, 90,68\% PyTorch). So với patch, cải thiện tương ứng 1,48
và 2,21 điểm; so với rows là 0,68 và 0,78 điểm. Các convolution học đặc trưng
cục bộ trước attention là giải thích kiến trúc hợp lý, nhưng chưa cô lập
tác động vì capacity và token features cũng thay đổi. Rows $ N=28$ có
attention matrix 784 phần tử/head; patch/stem $ N=49$ có 2.401, gấp 3,06 lần.
Chi phí toàn mạng còn gồm projections/FFN/CNN, nên không suy tốc độ chỉ từ $ N^2$.

\ReportFigure{assets/e2/head_ablation_curves.png}{Ablation H=2/4/8 trên manual CNN-stem, giữ D=64.}
H=2/4/8 đạt val 90,50/90,65/90,35\%; test 90,74/90,50/90,23\%.
H=4 tốt nhất theo validation nhưng H=2 cao hơn trên test; giữ quy tắc chọn
H=4 thay vì đổi quyết định sau test. Tăng head không tăng tham số projection,
mà chia feature mỗi head thành 32/16/8 chiều; accuracy không tăng đơn điệu.

\section{Learning curves, attention và lỗi}
\ReportFigure{assets/e2/learning_curves_loss.png}{Loss theo epoch, nhóm theo tokenizer; train nét đứt, validation nét liền.}
\ReportFigure{assets/e2/learning_curves_accuracy.png}{Accuracy theo epoch của các cấu hình E2.}
Các train--val gap khoảng 1,53--2,72 điểm phần trăm. Đọc curves cùng best
epoch cho thấy cần dừng khi validation chững, thay vì chọn loss train thấp nhất.
\ReportFigure{assets/e2/attention_best_manual.png}{Attention của manual CNN-stem H=4 trên hai ảnh validation cố định; trung bình head, block cuối, gộp 7 nhóm query/key, tổng mỗi hàng 100\%.}
Map thể hiện phân bố trọng số giữa các nhóm token sau CNN-stem, không phải
saliency trực tiếp ở pixel gốc. Việc trung bình head và gộp 7 nhóm làm mất
chi tiết từng head/token. Các ô sáng cho biết trọng số kết hợp lớn, không
tự chứng minh token là nguyên nhân của quyết định phân loại.

\begin{table}[H]\centering\caption{Nhầm lẫn lớn nhất của ba tokenizer manual H=4.}
\small\begin{tabular}{llrr}\toprule Model & Thật $\to$ dự đoán & Số ảnh & \% lớp thật\\\midrule
M-R4 & Shirt $\to $ T-shirt/top &139&13,9\\
M-P4 & Coat $\to $ Pullover &127&12,7\\
M-C4 & Shirt $\to $ T-shirt/top &94&9,4\\\bottomrule\end{tabular}\end{table}
\ReportFigure[.9]{assets/e2/failure_examples_test.png}{Gallery lỗi của các mô hình đại diện E2 chọn bằng validation.}
Ảnh lỗi và ma trận đầy đủ ở Phụ lục cho thấy nhóm áo vẫn khó tách.
Patch hay nhầm Coat/Pullover; stem giảm cặp Shirt/T-shirt/top trong bản
manual, song không xóa lỗi. E2 sử dụng hai cách cài đặt và ba cách tạo token;
giới hạn là encoder nhỏ, một seed, chưa kiểm soát capacity giữa tokenizers
và attention visualization chỉ mang tính mô tả.
